% !TEX program = xelatex
\documentclass[12pt,a4paper]{article}
\usepackage{fontspec}
\usepackage{polyglossia}
\setmainlanguage{russian}
\setotherlanguage{english}
\setmainfont{DejaVu Serif}
\setsansfont{DejaVu Sans}
\setmonofont{DejaVu Sans Mono}
\usepackage{geometry}
\geometry{margin=2.5cm}
\usepackage{amsmath,amssymb,amsthm,mathtools,bm}
\usepackage{microtype}
\usepackage{enumitem}
\usepackage{hyperref}
\hypersetup{colorlinks=true, linkcolor=blue, citecolor=blue, urlcolor=blue}
\newtheorem{definition}{Определение}
\newtheorem{axiom}{Аксиома}
\newtheorem{theorem}{Теорема}
\newtheorem{proposition}{Утверждение}
\newtheorem{remark}{Замечание}
\DeclareMathOperator{\Valid}{ValidInterface}
\DeclareMathOperator{\Interface}{Interface}
\DeclareMathOperator{\Auto}{Auto}
\DeclareMathOperator{\AlwaysAct}{AlwaysAct}
\DeclareMathOperator{\Input}{Input}
\DeclareMathOperator{\Adm}{Adm}
\DeclareMathOperator{\Dir}{Dir}
\DeclareMathOperator{\AutoAct}{AutoAct}
\DeclareMathOperator{\PassiveChannel}{PassiveChannel}
\DeclareMathOperator{\Control}{Control}
\DeclareMathOperator{\Der}{Der}
\DeclareMathOperator{\Struct}{Struct}
\DeclareMathOperator{\Dist}{Dist}
\DeclareMathOperator{\End}{End}
\DeclareMathOperator{\Tr}{Tr}
\DeclareMathOperator{\Attr}{Attr}
\DeclareMathOperator{\Coh}{Coh}
\DeclareMathOperator{\Sync}{Sync}
\DeclareMathOperator{\Loss}{Loss}
\DeclareMathOperator{\Stab}{Stab}
\newcommand{\M}{\mathcal M}
\newcommand{\B}{\mathcal B}
\newcommand{\Hcal}{\mathcal H}
\newcommand{\Ispace}{\mathfrak I}
\newcommand{\XiAuto}{\Xi_I^{\mathrm{auto}}}
\newcommand{\XiMan}{\mathcal M_{\Xi}^{\mathrm{auto}}}
\newcommand{\NXi}{\mathcal N_{\Xi}}
\newcommand{\PiAmp}{\Pi_{\mathrm{amp}}}
\newcommand{\PiStr}{\Pi_{\mathrm{str}}}
\newcommand{\PiStruct}{\Pi_{\mathrm{struct}}}

\title{\textbf{TMI-Core 0.2: Аксиоматическое ядро автоматических интерфейсов}\\
\large Первичный интерфейс, дефиниционная автоматичность и переход к автоматически структурированным полям}
\author{Salkutsan Aleksey}
\date{2026}
\begin{document}
\maketitle
\begin{abstract}
Данный документ фиксирует обновлённое ядро Теории многообразных интерфейсов: \(\mathcal T_{MI}^{0.2}\). Центральный переход версии 0.2 состоит в том, что автоматичность больше не рассматривается как внешняя надстройка над интерфейсом. Она вводится как дефиниционное свойство действительного интерфейса: если структура не способна самостоятельно интерпретировать допустимый вход и сформировать понятный выход для другой системы, она является не полноценным интерфейсом, а пассивным каналом или частью метаинтерфейса. Ядро связывает аксиому первичного интерфейса, дефиниционную автоматичность, различение интерфейса и пассивного канала, производные интерфейсы, проекцию на физический слой и масштабирование автоматически структурированных интерфейсных полей.
\end{abstract}
\tableofcontents
\newpage

\section{Назначение документа}
Цель этого документа - зафиксировать минимальное ядро, которое больше не расширяется новыми крупными сущностями. Дальнейшая работа должна уточнять определения, доказывать внутренние следствия и выводить проверяемые физические или инженерные проекции.

\section{Первичные объекты}
\begin{definition}[Нулевая недифференцированность]
\(N_0\) - предельное состояние без различия, сторон, отношения, порядка и направления. Это не физический вакуум и не объект природы, а предельный символ недифференцированности внутри аксиоматической схемы.
\end{definition}

\begin{definition}[Первичный интерфейс]
\(I_0\) - первая граница различения, через которую возникает минимальное различие.
\end{definition}

\begin{definition}[Первое различие]
\(\Delta_0\) - минимальный результат действия первичного интерфейса. Оно разворачивается в минимальную двустороннюю структуру \((A,\neg A)\).
\end{definition}

\section{Аксиома первичного интерфейса}
\begin{axiom}[Первичный интерфейс]
\begin{equation}
\boxed{A_{I_0}:\quad N_0\xrightarrow{I_0}\Delta_0\to(A,\neg A).}
\end{equation}
Краткая форма:
\begin{equation}
\boxed{N_0\xrightarrow{I_0}(A,\neg A).}
\end{equation}
\end{axiom}

Смысл аксиомы: первичный интерфейс не соединяет заранее готовые стороны, а создаёт саму возможность сторон. До \(I_0\) нет ни \(A\), ни \(\neg A\). После возникновения границы появляется различие \(\Delta_0\), а различие разворачивается в минимальную структуру \((A,\neg A)\).

\section{Действительный интерфейс}
\begin{definition}[Действительный интерфейс]
Действительный интерфейс \(I\) - структура, которая при наличии допустимого входа выполняет перевод, интерпретацию или структурирование между системами без внешнего управления каждым актом перевода.
\end{definition}

В рабочей форме:
\begin{equation}
I_{A\to B}\equiv\bigl(\partial(A,B),T_I(S_A\to S_B),\operatorname{Interp}_I,\Auto_I\bigr),
\end{equation}
где \(\partial(A,B)\) - граница взаимодействия, \(T_I\) - оператор перевода состояний, \(\operatorname{Interp}_I\) - операция интерпретации, а \(\Auto_I\) - самовыполнение интерфейсной функции при допустимом входе.

\section{Аксиома автоматичности интерфейса}
\begin{axiom}[Дефиниционная автоматичность]
\begin{equation}
\boxed{A_{\mathrm{auto}}:\quad \forall I\in\Ispace,\;\Valid(I)\Rightarrow\Auto(I).}
\end{equation}
Кратко:
\begin{equation}
\boxed{I\in\Ispace\Rightarrow\Auto(I).}
\end{equation}
\end{axiom}

Автоматичность не означает постоянную активность:
\begin{equation}
\boxed{\Auto(I)\not\Rightarrow\AlwaysAct(I).}
\end{equation}

Автоматичность означает самозапуск при наличии допустимого входа:
\begin{equation}
\boxed{\Auto(I)\equiv\bigl(\Input(I)\land\Adm(I)\land\Dir(I)\Rightarrow\AutoAct(I)\bigr).}
\end{equation}

\section{Пассивный канал и метаинтерфейс}
\begin{definition}[Пассивный канал]
Пассивный канал - структура, через которую может проходить сигнал или состояние, но которая не выполняет собственную интерпретационную функцию без внешнего управляющего агента.
\end{definition}

Если переводом управляет внешний агент \(U_{\mathrm{ext}}\), то действительным интерфейсом является составная структура:
\begin{equation}
\boxed{I'=(I,U_{\mathrm{ext}}).}
\end{equation}

\begin{equation}
\boxed{\PassiveChannel(I)\land\Control(U_{\mathrm{ext}},I)\Rightarrow\Interface(I').}
\end{equation}

\begin{theorem}[Различение интерфейса и канала]
Если структура не выполняет интерфейсную функцию автоматически при допустимом входе, то она либо не является действительным интерфейсом, либо является пассивным каналом, входящим в более широкий метаинтерфейс:
\begin{equation}
\boxed{\neg\Auto(I)\Rightarrow \neg\Valid(I)\lor\PassiveChannel(I).}
\end{equation}
\end{theorem}

\begin{proof}
По аксиоме автоматичности действительный интерфейс удовлетворяет \(\Valid(I)\Rightarrow\Auto(I)\). Контрапозиция даёт \(\neg\Auto(I)\Rightarrow\neg\Valid(I)\). Если при этом структура всё же участвует в передаче, но не выполняет перевод сама, она является пассивным каналом в составе более широкой системы \(I'=(I,U_{\mathrm{ext}})\). Следовательно, утверждение выполняется внутри аксиоматической схемы.
\end{proof}

\section{Производные интерфейсы}
\begin{definition}[Производный интерфейс]
\(\Der(I,I_0)\) означает, что интерфейс \(I\) является производным от первичного интерфейса \(I_0\), то есть принадлежит пространству интерфейсов, возникающему после первичного различения.
\end{definition}

\begin{equation}
\boxed{\Der(I,I_0)\Rightarrow I\in\Ispace\Rightarrow\Auto(I).}
\end{equation}

Для производных интерфейсов автоматичность раскрывается как условная:
\begin{equation}
\boxed{\Der(I,I_0)\Rightarrow\bigl(\Input(I)\land\Adm(I)\land\Dir(I)\Rightarrow\AutoAct(I)\bigr).}
\end{equation}

\section{Мост к физическому слою}
Если физическое поле является интерфейсным полем,
\begin{equation}
\Xi_I\in\Ispace_{\mathrm{phys}},
\end{equation}
то оно автоматично как интерфейс:
\begin{equation}
\boxed{\Xi_I\in\Ispace_{\mathrm{phys}}\Rightarrow\Auto(\Xi_I).}
\end{equation}
Следовательно, физическое интерфейсное поле должно рассматриваться как автоматически структурированное:
\begin{equation}
\boxed{\Xi_I\in\Ispace_{\mathrm{phys}}\Rightarrow\XiAuto.}
\end{equation}

\section{Проекционная цепочка}
Полная форма поля:
\begin{equation}
\boxed{\XiAuto.}
\end{equation}
Её проекции:
\begin{equation}
\boxed{\XiAuto\to\Xi_I\to\chi_I.}
\end{equation}
Здесь \(\chi_I\) - минимальная скалярная амплитудная проекция:
\begin{equation}
\boxed{\chi_I=\PiAmp(\XiAuto).}
\end{equation}

\section{Масштабирование}
Локальное автоматически структурированное поле масштабируется в сеть и многообразие:
\begin{equation}
\boxed{\XiAuto\to\{\Xi_{I,k}^{\mathrm{auto}}\}_{k\in K}\to\NXi\to\XiMan\to\Struct_I(\M).}
\end{equation}

\section{Объединённая формула ядра}
\begin{equation}
\boxed{N_0\xrightarrow{I_0}\Delta_0\to(A,\neg A)\Rightarrow\Ispace\Rightarrow\Auto(I).}
\end{equation}
\begin{equation}
\boxed{I_0\Rightarrow\Ispace\Rightarrow\Auto(I)\Rightarrow\XiAuto\Rightarrow\XiMan\Rightarrow\Struct_I(\M).}
\end{equation}

\section{Статус утверждений}
\begin{itemize}
\item Уровень A: аксиомы и определения, например \(\Valid(I)\Rightarrow\Auto(I)\).
\item Уровень B: внутренние теоремы, например переход от \(\XiAuto\) к \(\Struct_I\) при условиях стабилизации.
\item Уровень C: физические гипотезы, например существование \(\XiAuto\) как физического механизма.
\item Уровень D: проверяемые следствия, например увеличение времени когерентности в квантовом эксперименте.
\end{itemize}

\section{Заключение}
TMI-Core 0.2 фиксирует центральный стержень теории: первичный интерфейс создаёт возможность различия; действительный интерфейс автоматичен по определению; физические интерфейсные поля являются автоматически структурированными; множество таких полей образует масштабируемое интерфейсное многообразие, способное стабилизироваться как структура.

\begin{thebibliography}{9}
\bibitem{TMIcore} Salkutsan Aleksey. \textit{Theory of Manifold Interfaces. An Axiomatic Conceptual-Mathematical Framework of Interfaciality, Meaning, and Cognition}. Zenodo. \href{https://doi.org/10.5281/zenodo.19940675}{doi:10.5281/zenodo.19940675}.
\bibitem{TMIphysics} Salkutsan Aleksey. \textit{Theory of Manifold Interfaces: A Regularized Effective Physical Layer, Observable Consequences, and Limitations}. 2026.
\end{thebibliography}
\end{document}
